\chapter{Conclusion}
\label{cha:conclusion}
This study investigated the effectiveness of Logistic Regression, Random Forest, and Neural Network models for intrusion detection on the NSL-KDD dataset, addressing both binary and multi-class classification settings under severe class imbalance. The results show that, while all models achieve high accuracy in the binary task, performance drops substantially on the official test set, highlighting limited generalization to unseen attack patterns. In the multi-class setting, aggregate metrics hide important weaknesses in rare attack detection, highlighting the need for attack-focused evaluation. Logistic Regression over-generalizes, Random Forest favors dominant patterns, and Neural Networks improve minority recall but introduce higher complexity and class confusion. Overall, Random Forest provides the most balanced trade-off between performance, interpretability, computational efficiency, and stability, despite its limited effectiveness on rare attack detection. The primary insight of this work is that model choice alone is insufficient to address multi-class intrusion detection challenges. Future progress will depend on improved data representations, class-aware evaluation, and validation under more realistic conditions.